\section{Independent finite derivation of the inverse-quotient identities}\label{independent-finite-derivation-of-the-inverse-quotient-identities}

\subsection{Scope, inputs and reading record}\label{scope-inputs-and-reading-record}

This note verifies equations (3)--(31) of the user-supplied manuscript *The inverse-power theorem gives a quantitative reduction of the original observation complex*. It is an independent finite-dimensional derivation, not a claim to have re-proved the arithmetic inverse-power estimate, the Gamma division estimate, the monic-profile limit, or a Riemann-hypothesis endpoint.

The supplied text is at \nolinkurl{SUPPLIED_INVERSE_SECTOR_REDUCTION.md}. Its equations (3)--(31) were read in full. The preceding programme manuscript *The original observed inverse-power flag and its relative late heat*, \nolinkurl{INVERSE_FLAG_AND_RELATIVE_HEAT.tex}, was read through IFH10 to check the identity of the imported finite comparison. Its public source is \href{https://github.com/KokunoYumeto/zeta-function-research-reader/blob/edd4008690952fb8568a01744e49ad7e180401e9/workbenches/splitzero-tandem/continuations/20260921-original-observation/INVERSE_FLAG_AND_RELATIVE_HEAT.tex}{the complete IFH proof}. The exact source SHA-256 is \nolinkurl{468c92ee4224a3659be44f122efb9fbfc0c2a469624ff5da6d16280aa069a6ca}. In particular IFH1 defines the original observation metric, and IFH2--IFH7 give the finite source/observation sandwich and its passage to attained subquotients. This is programme mathematics, not an original human primary paper. No external primary-source theorem is silently attributed to it.

The mathematical task received from the root was to derive the Schur, quotient, chain-homotopy, inverse-division and observed-feedback identities, preserving the original objects and identifying any exact correction with a proof. Root retains source tracing and all later equations. No owner source, remote repository, or publication was changed by this task.

\subsection{1. Original finite objects and the minimum formula}\label{original-finite-objects-and-the-minimum-formula}

All adjoints denoted by a star in coordinate formulas are conjugate transposes. A metric matrix records a Hermitian form, not a replacement of the original form by the identity.

Let the full original polynomial be monic of degree q with Q(0) nonzero. The algebra is E=C{[}y{]}/(Q), represented in the complete ordered basis 1,y,\ldots,y\^{}(q−1). Multiplication by y is the invertible operator M. Let G be the specified positive Hermitian metric on E, let Λ:E→B be onto, and let I\_K:C\^{}m→E be the fixed full-column frame of K=ker Λ. Set

\[
H_K=I_K^*GI_K,\qquad
Q_B=(\Lambda G^{-1}\Lambda^*)^{-1},\qquad
L_B=G^{-1}\Lambda^*Q_B.
\]

For every onto matrix A from a space with metric H, the unique least-norm solution of Ax=b is

\[
S_Ab=H^{-1}A^*(AH^{-1}A^*)^{-1}b.
\]

Indeed \(AS_A=I\), and for \(z\in\ker A\), \(z^*HS_Ab=z^*A^*(AH^{-1}A^*)^{-1}b=0\). Every solution is \(S_Ab+z\); its squared norm is \(b^*(AH^{-1}A^*)^{-1}b+z^*Hz\). This proves attainment, uniqueness, the quotient metric, and the fact that S\_A is an isometry from that quotient metric. It also proves A is a contraction and that S\_AA is the original-metric orthogonal projection onto (ker A)\^{}⊥.

Apply this to Λ. Then P\_B=L\_BΛ is the original G-orthogonal projection onto K\^{}⊥. The complementary projection is

\[
P_K=I_KH_K^{-1}I_K^*G,\qquad P_B+P_K=I_E.
\]

For 0≤r≤q, retain the ordered inverse-power map

\[
U_rd=\sum_{j=1}^r d_j[y^{-j}],\qquad V_r=\operatorname{im}U_r.
\]

It is injective: multiplying U\_rd=0 by M\^{}r yields \(\sum_{j=1}^rd_jy^{r-j}=0\) in E, and the representative has degree less than q. Thus all its coefficients vanish. In the observation sector under consideration the imported IP lower factor is positive, so Φ\_r=ΛU\_r is also injective. Equivalently K∩V\_r=0. Define

\[
H^E=U_r^*GU_r,\quad H^B=\Phi_r^*Q_B\Phi_r,\quad
\Gamma=I_K^*GU_r.
\]

The finite input used below is precisely \(aH^E\preceq H^B\preceq H^E\), with a\textgreater0. Because H\^{}E is positive, this also entails a≤1. The stated arithmetic IP theorem supplies a and its lower bound exp(−ε); this finite derivation does not change that input.

\subsection{2. The two Schur complements, their eigenvalues and their exact loss}\label{the-two-schur-complements-their-eigenvalues-and-their-exact-loss}

Since P\_B is G-self-adjoint and idempotent,

\[
H^B=U_r^*GP_BU_r
=H^E-U_r^*GP_KU_r
=H^E-\Gamma^*H_K^{-1}\Gamma. \tag{F3}
\]

For \(z\in\mathbb C^m\) the quotient norm of I\_Kz modulo V\_r is the attained minimum of \((I_Kz+U_rd)^*G(I_Kz+U_rd)\) over all d.~Completing the exact quadratic form yields the minimizing coefficient \(d=-(H^E)^{-1}\Gamma^*z\) and hence

\[
\widehat H_K=H_K-\Gamma(H^E)^{-1}\Gamma^*. \tag{F4}
\]

This minimum remains positive because K∩V\_r=0. The minimum representative is \((I-P_V)I_Kz\), where \(P_V=U_r(H^E)^{-1}U_r^*G\) is the original orthogonal projection onto V\_r.

For the finite spectral comparison define, without altering either physical metric,

\[
T=H_K^{-1/2}\Gamma(H^E)^{-1/2}.
\]

The original metric matrices factor exactly as

\[
H^B=(H^E)^{1/2}(I_r-T^*T)(H^E)^{1/2},\qquad
\widehat H_K=H_K^{1/2}(I_m-TT^*)H_K^{1/2}.
\]

The positive eigenvalues of T*T and TT* coincide, with multiplicity: if \(T^*Tv=\mu v\) and \(\mu>0\), then \(Tv\ne0\) and \(TT^*(Tv)=\mu Tv\); the inverse correspondence on this eigenspace is \(\mu^{-1}T^*\). Their number is \(\nu=\operatorname{rank}T=\operatorname{rank}\Gamma\). The remaining eigenvalues are zero. IP gives \(0\preceq T^*T\preceq(1-a)I_r\), so the same nonzero-eigenvalue calculation gives \(0\preceq TT^*\preceq(1-a)I_m\). Consequently

\[
aH_K\preceq\widehat H_K\preceq H_K. \tag{F5}
\]

One can obtain the determinant equality without a spectral argument by taking the determinant of the same block Gram \(\begin{psmallmatrix}H_K&\Gamma\\\Gamma^*&H^E\end{psmallmatrix}\) in its two block orders. The equality is

\[
\det H_K\det H^B=\det H^E\det\widehat H_K,
\quad
\delta:=\log\frac{\det H^E}{\det H^B}
=\log\frac{\det H_K}{\det\widehat H_K}
=-\log\det(I_r-T^*T). \tag{F6}
\]

If the ν nonzero eigenvalues of T*T are μ\_j, then \(\delta=-\sum_{j=1}^{\nu}\log(1-\mu_j)\). Every factor 1−μ\_j lies in {[}a,1{]}. Thus

\[
0\le\delta\le\nu\log(1/a)\le\min(m,r)\varepsilon
\quad\text{when }a\ge e^{-\varepsilon}. \tag{F7}
\]

This proves the precise rank loss in (7), not a loss counted r times when r\textgreater m. With the manuscript's stated uniform ε=o(q) and m=O(k), it implies δ=o(kq) exactly as in (8). This is a propagation of the stated input, not a separate verification of its arithmetic range.

For the four-cutoff signs +,+,−,−, all four losses lie in the same interval {[}0,d{]}, where d=min(m,r)ε. Their signed sum belongs to {[}−2d,2d{]}: the sum of the two positively signed losses and the sum of the two negatively signed losses each belong to {[}0,2d{]}. Since log det Hhat=log det H\_K−δ, this proves the factor two in (9). No cancellation beyond the common nonnegative interval has been assumed.

\subsection{3. Explicit quotient coordinates and their attained metric}\label{explicit-quotient-coordinates-and-their-attained-metric}

Let rem\_Q denote the full degree-below-q representative. Define

\[
\pi_r[p]=\mathscr D^r(\operatorname{rem}_Q(y^rp)),\qquad
\mathscr D^rf=(f-\operatorname{Tay}_{<r}f)/y^r.
\tag{F10}
\]

This is well-defined on E, because replacing p by p+Qa does not change rem\_Q(y\^{}rp). It maps into P\_\textless q−r. Its kernel consists exactly of classes whose M\^{}r-image is in P\_\textless r. Since M\^{}rV\_r=P\_\textless r, invertibility of M\^{}r gives ker π\_r=V\_r. Let ι\_r be the ordinary degree-below-(q−r) inclusion. For p in this coefficient space, y\^{}rp has degree below q and its first r coefficients vanish; hence π\_rι\_rp=p.~Thus π\_r is onto.

The general minimum calculation in Section 1 yields

\[
\overline G_r=(\pi_rG^{-1}\pi_r^*)^{-1},\qquad
S_r=G^{-1}\pi_r^*\overline G_r. \tag{F11}
\]

Here S\_r is the minimum section, and S\_rπ\_r=I−P\_V. The quotient coordinate metric is therefore the attained one, not the raw polynomial Gram. Direct substitution gives

\[
(\pi_rI_K)^*\overline G_r(\pi_rI_K)
=I_K^*G(I-P_V)I_K=\widehat H_K.
\]

The physical substitution \(S=c+iy\) is retained by its full triangular map \[
(T_c)_{ab}=\binom ba c^{b-a}i^a\quad(a\le b),\qquad (T_c)_{ab}=0\quad(a>b).
\] It sends \(S\)-coefficients to \(y\)-coefficients. Therefore \(I_K^y=T_cI_K^S\) and \(G^S=T_c^*G^yT_c\), and both sides give the same fixed-kernel Gram. The diagonal phases \(i^a\) are retained.

\subsection{4. The quotient observation, its metric and the chain contraction}\label{the-quotient-observation-its-metric-and-the-chain-contraction}

Because Φ\_r is injective, choose r rows I with invertible minor Φ\_I, and let J be the remaining row indices. With coordinate selectors R\_I and R\_J, define the fixed algebraic quotient matrix

\[
\rho_r=R_J-\Phi_J\Phi_I^{-1}R_I.
\]

It kills Φ\_r. Its restriction to the coordinate complement supported on J is the identity, so it is onto and has kernel W\_r=im Φ\_r. This formula depends on the original algebraic maps, not on the cutoff metric. Define

\[
\overline\Lambda_r=\rho_r\Lambda\iota_r.
\]

Since x−ι\_rπ\_rx lies in V\_r and ρ\_rΛ kills V\_r, barΛ\_rπ\_r=ρ\_rΛ. Onto-ness follows from that of ρ\_rΛ. If barΛ\_rπ\_rx=0, then Λx=Λv for some v∈V\_r, so x−v∈K and π\_rx∈π\_rK. Conversely π\_rK is killed. The restriction π\_r:K→ker barΛ\_r is injective and onto. This proves all arrows and exactness in (12)--(13).

For the target quotient the attained metric is

\[
\overline Q_r=(\rho_rQ_B^{-1}\rho_r^*)^{-1}.
\]

The inverse of this matrix is exactly \[
\begin{aligned}
\rho_r\Lambda G^{-1}\Lambda^*\rho_r^*
&=\overline\Lambda_r\pi_rG^{-1}\pi_r^*\overline\Lambda_r^*\\
&=\overline\Lambda_r\,\overline G_r^{-1}\overline\Lambda_r^* .
\end{aligned}
\] Inverting proves (14) and equality of the two orders of complete minimum, without identifying either quotient metric with a raw coefficient Gram.

Now set

\[
h=U_r(H^B)^{-1}\Phi_r^*Q_B:B\longrightarrow E. \tag{F15}
\]

Then hΦ\_r=U\_r, so hΛ is idempotent with image V\_r. Its kernel is the preimage under Λ of the Q\_B-orthogonal complement to W\_r. Meanwhile \(P_W=\Lambda h=\Phi_r(H^B)^{-1}\Phi_r^*Q_B\) is the original Q\_B-orthogonal projection onto W\_r. The degree-zero projection hΛ need not be orthogonal.

Write hb=U\_rd. The equality d=(H\^{}B)\^{}(−1)Φ\_r*Q\_Bb gives

\[
\|hb\|_G^2=d^*H^Ed\le a^{-1}d^*H^Bd
=a^{-1}\|P_Wb\|_{Q_B}^2\le a^{-1}\|b\|_{Q_B}^2.
\]

Thus \(\|h\|\le a^{-1/2}\), proving (16) with its exact metric domains.

Let \(S_\rho=Q_B^{-1}\rho_r^*\overline Q_r\) be the minimum section of ρ\_r. Since S\_ρρ\_r=I−P\_W, the fully explicit chain section of p=(π\_r,ρ\_r) is

\[
j^0=(I_E-h\Lambda)S_r,
\qquad j^1=S_\rho. \tag{F17a}
\]

First π\_rj\^{}0=π\_rS\_r=I and ρ\_rj\^{}1=I. Secondly,

\[
\Lambda j^0=(I_B-\Lambda h)\Lambda S_r
=S_\rho\rho_r\Lambda S_r
=S_\rho\overline\Lambda_r=j^1\overline\Lambda_r.
\]

Thus j is a chain map. The difference x−S\_rπ\_rx belongs to V\_r and I−hΛ kills V\_r; hence j\^{}0π\_r=I−hΛ. The other degree gives j\^{}1ρ\_r=I−Λh. Regard h:B→E as a degree −1 map, with its other component zero. For the differential d=Λ these two equalities are exactly

\[
I-jp=dh+hd. \tag{F17b}
\]

This proves the finite chain-homotopy equivalence (18). It fixes K pointwise: for x∈K, j\^{}0π\_rx=(I−hΛ)x=x. No infinite theta complex is introduced by this calculation.

\subsubsection{Complete polynomial relation spaces remain present}\label{complete-polynomial-relation-spaces-remain-present}

Let the actual source P\_N have its positive source metric, and let J\_N:P\_N→E be the full remainder map, with minimum section L\_N. Then J\_NL\_N=I and L\_N is an isometry from G. The lifted homotopy \(\widetilde h=L_Nh\) has exactly the same norm. Its range is L\_NV\_r. Moreover \(\widetilde h\Lambda J_N\) vanishes on the complete original relation space \(\ker J_N=Q\mathcal P_{N-q}\), and this relation space is orthogonal to L\_NE and therefore to L\_NV\_r.

There is a corresponding reduction of {[}P\_N→B{]} by the contractible subcomplex {[}L\_NV\_r→W\_r{]}. Its source quotient is P\_N/(L\_NV\_r), not an identification of P\_N with E: every original Q-relation remains in degree-zero cohomology. The same proof with differential ΛJ\_N and homotopy L\_Nh gives the contraction. This specifies exactly what the manuscript's polynomial-source lift preserves.

\subsection{5. Exact determinant allocation and its fixed frame factor}\label{exact-determinant-allocation-and-its-fixed-frame-factor}

The Gram of the augmented kernel K+V\_r in the fixed ordered frame {[}I\_K,U\_r{]} has blocks H\_K,Γ,Γ*,H\^{}E. Its determinant is

\[
\det([I_K,U_r]^*G[I_K,U_r])=\det H_K\det H^B. \tag{F19}
\]

Choose any fixed right inverse Z of ρ\_r and set C={[}Φ\_r,Z{]}. This square matrix is invertible. In that frame the Q\_B Gram has leading block H\^{}B and lower Schur complement barQ\_r. Therefore the complete equality is

\[
|\det C|^2\det Q_B=\det H^B\det\overline Q_r,
\]

or

\[
\log\det\overline Q_r
=\log\det Q_B-\log\det H^B+2\log|\det C|. \tag{F20}
\]

C is fixed across the four cutoffs, so the last term cancels under their signs. This is precisely (20), including the coefficient-frame constant before the return is taken. Combined with (F19), the signs of the two redistributed terms in (21) are + and −. Any use of the evaluated coefficient Cfrak belongs to the imported IFH coefficient theorem; the allocation itself has no asymptotic error.

\subsection{6. Inverse multiplication as exact Taylor division between quotients}\label{inverse-multiplication-as-exact-taylor-division-between-quotients}

For x∈E retain its actual minimum source representative L\_Nx, and write

\[
L_Nx=\sum_{j=0}^{s-1}c_j(x)y^j+y^s\mathscr D^sL_Nx.
\]

Apply J\_N, and then multiply by M\^{}(−s). The result is

\[
M^{-s}x=J_N\mathscr D^sL_Nx+
\sum_{n=1}^sc_{s-n}(x)[y^{-n}]. \tag{F22}
\]

Thus the ordering in F\_(s,N) is exactly (F\_(s,N)x)\emph{n=c}(s−n)(x), as in (22). No assertion has been made that polynomial Taylor coefficients themselves are well-defined on E; the map L\_N supplies their specific representative.

For the norm in (23), the supplied Gamma comparison has the form ell\_(k,N) H\_Gamma≤H\_native on degrees ≤N and H\_native≤u\_k H\_Gamma on the output range. Apply these two comparisons on the outside of the s-fold division estimate, giving

\[
\|\mathscr D^sf\|_{\rm native}
\le \sqrt{u_k}\|\mathscr D^sf\|_\Gamma
\le\sqrt{u_k}(a_N^\Gamma)^s\|f\|_\Gamma
\le\sqrt{u_k/\ell_{k,N}}(a_N^\Gamma)^s\|f\|_{\rm native}.
\]

The source minimum L\_N is an isometry and J\_N is a contraction, so \textbar\textbar J\_N D\^{}s L\_N\textbar\textbar≤B\_(s,N). Only one comparison ratio occurs. This derives the use of the imported division bound, not its original Gamma proof.

Since \(M^{-s}V_r\subset V_{r+s}\), the map

\[
\mathcal I_{r,s}=\pi_{r+s}M^{-s}\iota_r:E/V_r\to E/V_{r+s}
\]

is well-defined for r+s≤q. For the canonical coordinate p of degree below q−r, calculate

\[
\pi_{r+s}M^{-s}\iota_rp
=\mathscr D^{r+s}\operatorname{rem}_Q(M^{r+s}M^{-s}p)
=\mathscr D^{r+s}(y^rp)=\mathscr D^sp. \tag{F24}
\]

The second equality means the representative of the algebraic product: \(M^{r+s}M^{-s}p=M^rp\), whose representative is y\^{}rp since its degree is below q. Hence the kernel is exactly P\_\textless s inside P\_\textless q−r and the map is onto P\_\textless q−r−s. Equivalently, \(M^sV_{r+s}=V_r\oplus\mathcal P_{<s}\): multiplying by M\^{}r sends these two summands to P\_\textless r and \(\operatorname{span}(y^r,\ldots,y^{r+s-1})\), disjoint complete monomial ranges of degree below q. This proves (25), including its first injection.

For the norm, take the actual minimum representative x=S\_rz of a quotient class z. The last term of (F22) belongs to V\_s⊂V\_(r+s), so

\[
\mathcal I_{r,s}z=\pi_{r+s}J_N\mathscr D^sL_NS_rz.
\]

Each outer quotient map is a contraction and S\_r an isometry in the attained metrics. This gives (26). Directly from \(M^{-t}M^{-s}=M^{-s-t}\), or from \(\mathscr D^t\mathscr D^s=\mathscr D^{s+t}\), one obtains (27) with its degree-changing domain and codomain. No fixed quotient is being treated as M-invariant.

\subsection{7. Complete observed leakage and composition memory}\label{complete-observed-leakage-and-composition-memory}

Let barL\_r be the attained minimum section of barΛ\_r and let barP\^{}K\_r=I−barL\_r barΛ\_r. It is the original quotient-metric orthogonal projection onto ker barΛ\_r. The observed map is

\[
\mathcal B_{r,s}=\overline\Lambda_{r+s}
\mathcal I_{r,s}\overline L_r.
\]

The section is an isometry and barΛ\_(r+s) is a contraction. Consequently \textbar\textbar B\_(r,s)\textbar\textbar≤B\_(s,N), exactly (28). Insert I=barL\_r barΛ\_r+barP\^{}K\_r to obtain

\[
\overline\Lambda_{r+s}\mathcal I_{r,s}
=\mathcal B_{r,s}\overline\Lambda_r
+\overline\Lambda_{r+s}\mathcal I_{r,s}\overline P^K_r. \tag{F29}
\]

All three factors in the last term have norm ≤1, B\_(s,N), ≤1 in their actual metrics, proving the same leakage bound.

For the composition use the exact identity (27) and subtract the observed product. The middle difference is \(I-\overline L_{r+s}\overline\Lambda_{r+s}=\overline P^K_{r+s}\), so

\[
\begin{aligned}
\mathcal B_{r,s+t}-\mathcal B_{r+s,t}\mathcal B_{r,s}
={}&\overline\Lambda_{r+s+t}\mathcal I_{r+s,t}
\overline P^K_{r+s}\mathcal I_{r,s}\overline L_r.
\end{aligned}\tag{F30}
\]

The sign is positive. The norm is at most B\_(t,N)B\_(s,N). On the retained injective conductor domain ker barΛ\_(r+s)=π\_(r+s)K has dimension m, and the rank is therefore at most m. Even outside injectivity of a particular inverse sector, the kernel dimension is at most m because its kernel remains π\_(r+s)K; the formulas involving its quotient metric still use its actual target quotient dimension.

Finally ρ\_(r+s)ΛU\_s=0 since ΛV\_s⊂W\_(r+s). Apply ρ\_(r+s)Λ to (F22) and restrict the input by I\_K. This proves

\[
\rho_{r+s}\Lambda M^{-s}I_K
=\rho_{r+s}\Lambda J_N\mathscr D^sL_NI_K. \tag{F31}
\]

The frame I\_K is an isometry from its retained coordinate metric H\_K to its physical image, and ρ\_(r+s)Λ is a contraction for the attained target metric. This proves the stated norm bound, without quotienting or changing the metric on the original K.

\subsection{8. Exact finite verification and conclusion}\label{exact-finite-verification-and-conclusion}

The companion \texttt{verify\_exact.py} uses SymPy exact rational complex arithmetic, the full polynomial \(Q=y^6+2y^4-y^3+3y+2\), an eight-dimensional positive source Gram with nonzero imaginary off-diagonal entries, its original quotient Gram, and a fixed onto observation and fixed two-dimensional kernel frame. The source metric is constructed explicitly, then G and all sections are obtained by their stated attained-minimum formulas. No floating or identity-metric replacement is used.

The script checks r=1,2,3, full Schur complements, all quotient and chain maps, both homotopy identities, the exact frame determinant factor, inverse Taylor decomposition, and the observed feedback identity. It requires the observed composition defect to be nonzero, so success cannot result merely from accidentally selecting an invariant observed sector. The actual result count and script hash are in \nolinkurl{EXACT_CHECKS.json}.

No correction to the finite identities (3)--(31) was found. The two details made explicit here are already consistent with the supplied manuscript: (i) the fixed frame factor 2 log\textbar det C\textbar{} in (F20) exists before the four-cutoff return; (ii) the polynomial-source contraction quotients by L\_NV\_r and retains its entire Q-relation kernel, rather than replacing that larger source complex by {[}E→B{]}.
